\section{Even cyclic Clifford constraints and cross-rank consequences}
\label{sec:even-clifford-retry}

Both rank-six questions remain open. We extend the stabilizer-rotation
criterion to even orders, and identify precisely two consequences of the
existing structural reductions.

\subsection{Even cyclic characters}
Use the hypotheses of Proposition~\ref{prop:stabilizer-rotation}: an
order-two invertible $h$ admits the algebra $1+h$, a simple $Y$ is fixed
on both sides, and its normalized left and right $h$-actions anticommute.
On $V=\operatorname{Hom}(FZ,Y^{\otimes n})$, normalize generalized
rotation to an operator $E_0$ with $E_0^n=1$. The pair-insertion operators
$A_j$, indexed cyclically, satisfy
\[
 A_j^2=1,\qquad A_jA_{j+1}=-A_{j+1}A_j,
 \qquad A_jA_k=A_kA_j\quad(k\ne j\pm1,j),
 \qquad E_0A_jE_0^{-1}=A_{j+1}.
\]
The boundary relations follow from the same rotation naturality as in the
odd-order proof; normalization of $E_0$ does not change conjugation.

\begin{proposition}[Even cyclic character criterion]
Let $n=2m\ge4$ and $r=2^{m-1}$. The eigenvalue multiplicity vector of
$E_0$ is a nonnegative integral sum of cyclic shifts of the following
three multisets of exponents modulo $n$:
\begin{align*}
 &\left\{-\frac12\sum_{a\in S}a:
       S\subset\{1,3,\ldots,2m-1\},\ |S|\text{ even}\right\},\\
 &\left\{\frac{1-\sum_{a\in S}a}{2}:
       S\subset\{1,3,\ldots,2m-1\},\ |S|\text{ odd}\right\},\\
 &\bigcup_{S\subset\{1,\ldots,m-1\}}
       \left\{-\sum_{a\in S}a,\ m-\sum_{a\in S}a\right\}.
\end{align*}
Their dimensions are respectively $r,r,2r$. The divisions by two are
ordinary integer divisions of even numerators.
\end{proposition}
\begin{proof}
The adjacency pairing of the even cycle over $\mathbb F_2$ has
two-dimensional radical, generated by its alternating vectors. The
insertion algebra is consequently a direct sum of four copies of
$M_r(\mathbb C)$. Rotation fixes two blocks and exchanges two.

For the fixed blocks, represent the generators by adjacent Clifford
bilinears with antiperiodic boundary. The rotation has eigenangles
$(2a+1)\pi/n$, $0\le a<m$. Its two half-spin representations have
weights indexed by even and odd subsets, respectively. Rescaling their
implementers to have $n$th power one gives the first two multisets.
Any other rescaling adds a constant modulo $n$ to every exponent.

For the exchanged blocks, use periodic boundary. The orthogonal
transformation has angles $2\pi a/n$, $1\le a<m$, and eigenvalues
$1,-1$. A pin implementer has opposite eigenvalue pairs; after the same
normalization these give the third multiset. Its shifts have only $m$
distinct characters because shifting by $m$ preserves it.

Each fixed block has $n$ possible scalar implementers. The exchanged
pair has $m$ irreducibles induced from its rotation-square stabilizer.
Their squared dimensions sum to
\[
 2nr^2+m(2r)^2=n2^n,
\]
the dimension of the insertion algebra extended by cyclic rotation.
Thus these representations exhaust its irreducibles. Decomposing $V$
proves the assertion.
\end{proof}

For $R6\_4\_155$, the frozen generalized-rotation certificate supplies
orders $4,6,\ldots,24$, both noninvertible objects, and all $52$ central
targets. Every one of the $1144$ resulting multiplicity vectors has an
exact nonnegative integral character decomposition. The independent
integer replay also rejects a mutated witness. Hence the criterion
gives no exclusion or associator construction for this ring.

\subsection{Pointed automorphisms and the rank-eight alternative}
The earlier central-lift theorem for $R6\_4\_155$ distinguishes the four
pointed simples by their unique central twists $1,i,-1,-i$.
A tensor autoequivalence preserves the unique pivotal structure and
these lifts, and therefore fixes every pointed simple class. The
vanishing and uniqueness assertions transport through the established
Galois reduction, so this conclusion applies to arbitrary complex
categorifications. An invertible normalizer in an extension consequently
centralizes the Klein-four pointed subgroup. In particular it cannot
give the specified $D_8$ pointed extension. This excludes a rank-nine
Morita branch, not $R6\_4\_155$ itself.

The surviving rank-eight alternative in the established
$R6\_8\_1660$ Morita list is exactly $R8\_3\_941$. The basis permutation
from that alternative to the census basis is
\[
 (0,3,6,7,4,5,2,1).
\]
All $512$ structure constants agree. The statement remains conditional:
the nonpointed algebra producing this module category has not been
forced to exist. The guaranteed pointed algebra on the rank-eight side
does not supply its inverse.

Complete proofs, pinned inputs, exact character decompositions and the
cross-rank identification are in the rank-six folders of
\path{evidence/unresolved_retry_20261024/}. The numerical laptop search
is an incomplete heuristic campaign and supplies neither a model nor
an exclusion. No new rank-six classification decision is asserted.
